\documentclass[12pt,a4paper]{article}

% ------------------------------------------------
% PACOTES
% ------------------------------------------------
\usepackage[utf8]{inputenc}
\usepackage[T1]{fontenc}
\usepackage[brazil]{babel}
\usepackage{lmodern}
\usepackage{geometry}
\usepackage{amsmath,amssymb}
\usepackage{tikz}
\usetikzlibrary{automata,positioning,arrows.meta}
\usepackage{listings}
\usepackage{xcolor}
\usepackage{booktabs}
\usepackage{array}
\usepackage{float}
\usepackage{hyperref}
\usepackage{caption}

\geometry{margin=2.5cm}
\hypersetup{
    colorlinks=true,
    linkcolor=blue,
    urlcolor=blue,
    citecolor=blue
}

% ------------------------------------------------
% CONFIGURAÇÃO DOS DIAGRAMAS
% ------------------------------------------------
\definecolor{azul}{RGB}{35,90,150}
\definecolor{verde}{RGB}{35,125,90}
\definecolor{cinza}{RGB}{90,90,90}

\tikzset{
    afd/.style={
        ->,
        >=Stealth,
        auto,
        semithick,
        every state/.style={
            draw=azul,
            fill=blue!4,
            text=black,
            minimum size=10mm
        },
        accepting/.style={
            draw=verde,
            fill=green!5,
            double,
            double distance=1.5pt
        }
    }
}

% ------------------------------------------------
% CONFIGURAÇÃO DOS CÓDIGOS
% ------------------------------------------------
\lstdefinestyle{pythonstyle}{
    language=Python,
    basicstyle=\ttfamily\small,
    keywordstyle=\color{blue}\bfseries,
    commentstyle=\color{gray}\itshape,
    stringstyle=\color{verde},
    numbers=left,
    numberstyle=\tiny\color{gray},
    frame=single,
    breaklines=true,
    showstringspaces=false,
    tabsize=4,
    inputencoding=utf8,
    extendedchars=true
}
\lstset{style=pythonstyle}

% ------------------------------------------------
% IDENTIFICAÇÃO
% ------------------------------------------------
\title{\textbf{Analisador Léxico: Construção de um Scanner para a Linguagem MiniLang}}
\author{
Samuel Douglas \and Ryan Rodrigues \and Rafael Alexandre \and
Alânderson Santos \and Daniel Manuel\\[0.4em]
\small Ciência da Computação
}
\date{\today}

\begin{document}
\maketitle

\begin{abstract}
A análise léxica é a primeira etapa do processo de compilação e tem como
objetivo transformar o código-fonte em uma sequência de unidades denominadas
tokens. Este trabalho apresenta a especificação de um analisador léxico
(scanner) para uma linguagem fictícia denominada MiniLang. O estudo descreve
as categorias léxicas reconhecidas, a representação formal de Autômatos
Finitos Determinísticos (AFDs), seus diagramas de transição, tabelas de
transição e o processo de minimização. Também é apresentada uma implementação
de apoio em Python, utilizando expressões regulares para identificar
palavras-chave, identificadores, números inteiros, números decimais,
operadores e delimitadores. O objetivo é demonstrar como os conceitos de
Linguagens Formais e Autômatos podem ser aplicados à construção de um
componente inicial de um compilador.
\end{abstract}

\noindent\textbf{Palavras-chave:} Análise léxica; Autômato Finito Determinístico;
Compiladores; Scanner; MiniLang.

\section{Introdução}

Um compilador realiza diferentes etapas para transformar um programa escrito
em uma linguagem de programação em uma representação que possa ser processada
ou executada por um computador. A primeira dessas etapas é a análise léxica,
responsável por percorrer o código-fonte e agrupar sequências de caracteres
em unidades significativas chamadas \textit{tokens}.

O componente responsável por essa atividade é denominado analisador léxico
ou \textit{scanner}. Durante a leitura, ele identifica padrões como nomes de
variáveis, palavras reservadas, valores numéricos e símbolos operadores.
Essas unidades são encaminhadas às etapas posteriores do processo de
compilação.

Neste projeto, foi definida uma linguagem fictícia chamada MiniLang. O
analisador proposto reconhece um conjunto reduzido de elementos léxicos,
permitindo demonstrar o funcionamento dos autômatos e a classificação dos
tokens.

\subsection{Objetivo}

O objetivo geral é especificar e implementar um analisador léxico para a
MiniLang. Como objetivos específicos, busca-se:
\begin{itemize}
    \item apresentar as categorias de tokens reconhecidas;
    \item formalizar os autômatos por meio de quíntuplas;
    \item representar os estados e as transições utilizando diagramas TikZ;
    \item apresentar tabelas de transição e discutir a minimização dos AFDs;
    \item implementar uma versão de apoio em Python.
\end{itemize}

\section{Especificação léxica da MiniLang}

A MiniLang utiliza um conjunto de palavras-chave, identificadores, literais
numéricos, operadores e delimitadores. As categorias consideradas neste
trabalho são apresentadas na Tabela~\ref{tab:tokens}.

\begin{table}[H]
\centering
\caption{Categorias léxicas da MiniLang}
\label{tab:tokens}
\begin{tabular}{p{3.4cm}p{9.0cm}}
\toprule
\textbf{Categoria} & \textbf{Exemplos} \\
\midrule
Palavras-chave & \texttt{if}, \texttt{else}, \texttt{while}, \texttt{int}, \texttt{float} \\
Identificadores & \texttt{contador}, \texttt{idade}, \texttt{valor} \\
Números inteiros & \texttt{10}, \texttt{25}, \texttt{100} \\
Números decimais & \texttt{10.5}, \texttt{3.14} \\
Operadores aritméticos & \texttt{+}, \texttt{-}, \texttt{*}, \texttt{/} \\
Atribuição & \texttt{=} \\
Delimitadores & \texttt{(}, \texttt{)}, \texttt{\{}, \texttt{\}}, \texttt{;} \\
\bottomrule
\end{tabular}
\end{table}

Um identificador deve começar com uma letra ou sublinhado e pode ser seguido
por letras, dígitos ou sublinhados. Essa regra é representada pela expressão
regular:
\[
    [a-zA-Z\_][a-zA-Z0-9\_]^*
\]

Os números inteiros são formados por um ou mais dígitos:
\[
    [0-9]^+
\]

Os números decimais possuem uma parte inteira, um ponto e uma parte fracionária,
todas com pelo menos um dígito:
\[
    \texttt{[0-9]+[.][0-9]+}
\]

\section{Fundamentação formal: Autômato Finito Determinístico}

Um Autômato Finito Determinístico (AFD) é uma máquina abstrata utilizada para
reconhecer linguagens regulares. Formalmente, é definido pela quíntupla:
\begin{equation}
    M = (Q,\Sigma,\delta,q_0,F)
    \label{eq:afd}
\end{equation}

Em que:
\begin{itemize}
    \item $Q$ é o conjunto finito de estados;
    \item $\Sigma$ é o alfabeto de entrada;
    \item $\delta: Q \times \Sigma \rightarrow Q$ é a função de transição;
    \item $q_0 \in Q$ é o estado inicial;
    \item $F \subseteq Q$ é o conjunto de estados de aceitação.
\end{itemize}

No analisador léxico, os AFDs representam os padrões que podem ser
reconhecidos durante a leitura dos caracteres. Quando a entrada termina em
um estado de aceitação, o padrão correspondente foi reconhecido. Transições
não permitidas conduzem à rejeição ou a um estado de erro.

\section{Diagramas e especificação dos autômatos}

\subsection{AFD para identificadores}

O autômato de identificadores exige que o primeiro caractere seja uma letra
($L$) ou sublinhado (\texttt{\_}). Depois do primeiro caractere, são
permitidas letras, dígitos ($D$) e sublinhados. O estado $q_1$ é final,
pois representa uma sequência que atende à regra de identificador.

A expressão regular correspondente é:
\[
    [a-zA-Z\_][a-zA-Z0-9\_]^*
\]

A definição formal simplificada é:
\[
\begin{aligned}
Q &= \{q_0,q_1,q_d\},\\
\Sigma &= \{L,D,\_\},\\
q_0 &= \text{estado inicial},\\
F &= \{q_1\}.
\end{aligned}
\]
O estado $q_d$ representa o estado de erro. O diagrama é apresentado na
Figura~\ref{fig:afd-identificador}.

\begin{figure}[H]
\centering
\begin{tikzpicture}[afd,node distance=3.2cm]
    \node[state,initial] (q0) {$q_0$};
    \node[state,accepting,right=of q0] (q1) {$q_1$};
    \path
        (q0) edge node[above] {$L,\ \_$} (q1)
        (q1) edge[loop above] node {$L,\ D,\ \_$} (q1);
\end{tikzpicture}
\caption{AFD para reconhecimento de identificadores.}
\label{fig:afd-identificador}
\end{figure}

\subsection{AFD para números inteiros e decimais}

O reconhecimento de números utiliza cinco estados: $q_0$ é inicial; $q_1$
reconhece a parte inteira; $q_2$ indica que o ponto decimal foi lido; $q_3$
reconhece a parte fracionária; e $q_d$ representa erro.

Os estados finais são $q_1$ e $q_3$, pois uma sequência formada apenas por
dígitos é um inteiro, enquanto uma sequência com parte fracionária completa
é um decimal. A definição é:
\[
\begin{aligned}
Q &= \{q_0,q_1,q_2,q_3,q_d\},\\
\Sigma &= \{D,.\},\\
F &= \{q_1,q_3\}.
\end{aligned}
\]

A Figura~\ref{fig:afd-numeros} apresenta as transições principais. Uma
transição não representada leva ao estado de erro.

\begin{figure}[H]
\centering
\begin{tikzpicture}[afd,node distance=2.2cm]
    \node[state,initial] (q0) {$q_0$};
    \node[state,accepting,right=of q0] (q1) {$q_1$};
    \node[state,right=of q1] (q2) {$q_2$};
    \node[state,accepting,right=of q2] (q3) {$q_3$};
    \path
        (q0) edge node[above] {$D$} (q1)
        (q1) edge[loop above] node {$D$} (q1)
        (q1) edge node[above] {$.$} (q2)
        (q2) edge node[above] {$D$} (q3)
        (q3) edge[loop above] node {$D$} (q3);
\end{tikzpicture}
\caption{AFD para reconhecimento de números inteiros e decimais.}
\label{fig:afd-numeros}
\end{figure}

\subsection{AFD para operadores e delimitadores}

Os operadores e delimitadores são reconhecidos a partir de um único símbolo.
Para fins de classificação, são separados em três grupos: operadores
aritméticos, símbolo de atribuição e delimitadores. A definição de estados é:
\[
\begin{aligned}
Q &= \{q_0,q_1,q_2,q_3,q_d\},\\
\Sigma &= \{+,-,*,/,=,(,),\{,\},;\},\\
F &= \{q_1,q_2,q_3\}.
\end{aligned}
\]

O estado $q_1$ representa operadores aritméticos; $q_2$, atribuição; e
$q_3$, delimitadores. O diagrama está representado na
Figura~\ref{fig:afd-operadores}.

\begin{figure}[H]
\centering
\begin{tikzpicture}[afd]
    \node[state,initial] (q0) at (0,0) {$q_0$};
    \node[state,accepting] (q1) at (4,2) {$q_1$};
    \node[state,accepting] (q2) at (4,0) {$q_2$};
    \node[state,accepting] (q3) at (4,-2) {$q_3$};
    \path
        (q0) edge node[above,sloped] {$+,-,*,/$} (q1)
        (q0) edge node[above] {$=$} (q2)
        (q0) edge node[below,sloped] {$(,),\{,\},;$} (q3);
    \node[right=5mm of q1,align=left] {\small Operador\\[-1mm]\small aritmético};
    \node[right=5mm of q2] {\small Atribuição};
    \node[right=5mm of q3] {\small Delimitador};
\end{tikzpicture}
\caption{AFD para operadores e delimitadores.}
\label{fig:afd-operadores}
\end{figure}

Os estados finais são mantidos separados porque cada um corresponde a uma
categoria de token diferente. Se o objetivo fosse apenas reconhecer a união
dos símbolos, os estados finais poderiam ser equivalentes quanto à aceitação
da linguagem; entretanto, o analisador precisa preservar a classificação
produzida.

\section{Tabelas de transição}

As tabelas de transição descrevem o estado de destino para cada combinação
de estado atual e símbolo de entrada. Na Tabela~\ref{tab:trans-ident}, $L$
representa uma letra, $D$ representa um dígito e \texttt{\_} representa o
sublinhado.

\begin{table}[H]
\centering
\caption{Tabela de transição do AFD de identificadores}
\label{tab:trans-ident}
\begin{tabular}{c|ccc}
\toprule
$\delta$ & $L$ & $D$ & \texttt{\_} \\
\midrule
$\rightarrow q_0$ & $q_1$ & $q_d$ & $q_1$ \\
$*q_1$ & $q_1$ & $q_1$ & $q_1$ \\
$q_d$ & $q_d$ & $q_d$ & $q_d$ \\
\bottomrule
\end{tabular}
\end{table}

A Tabela~\ref{tab:trans-num} apresenta as transições do autômato de números.
O símbolo $*$ indica estado de aceitação e $\rightarrow$ indica estado
inicial.

\begin{table}[H]
\centering
\caption{Tabela de transição do AFD de números}
\label{tab:trans-num}
\begin{tabular}{c|cc}
\toprule
$\delta$ & $D$ & $.$ \\
\midrule
$\rightarrow q_0$ & $q_1$ & $q_d$ \\
$*q_1$ & $q_1$ & $q_2$ \\
$q_2$ & $q_3$ & $q_d$ \\
$*q_3$ & $q_3$ & $q_d$ \\
$q_d$ & $q_d$ & $q_d$ \\
\bottomrule
\end{tabular}
\end{table}

\section{Minimização dos AFDs}

A minimização busca obter um autômato equivalente com o menor número
possível de estados. Um método utilizado é o refinamento de partições:
inicialmente, os estados são separados entre finais e não finais. Em seguida,
os grupos são divididos quando seus estados apresentam transições para
partições diferentes.

\subsection{Identificadores}

Para o autômato de identificadores, a partição inicial é:
\[
    P_0 = \big\{\{q_1\},\{q_0,q_d\}\big\}.
\]
Os estados $q_0$ e $q_d$ são distinguíveis, pois, ao receber uma letra,
$q_0$ transita para o estado final $q_1$, enquanto $q_d$ permanece no estado
de erro. Assim, obtêm-se três classes:
\[
    P = \big\{\{q_0\},\{q_1\},\{q_d\}\big\}.
\]
Portanto, o AFD de identificadores já é mínimo.

\subsection{Números}

Para o autômato de números, a partição inicial é:
\[
    P_0 = \big\{\{q_1,q_3\},\{q_0,q_2,q_d\}\big\}.
\]
Os estados são separados porque apresentam comportamentos diferentes diante
dos símbolos de entrada. Por exemplo, com o símbolo ponto, $q_1$ transita
para $q_2$, enquanto $q_3$ transita para $q_d$. Com isso, $q_1$ e $q_3$ não
são equivalentes. O refinamento distingue também $q_0$, $q_2$ e $q_d$,
resultando em cinco classes para cinco estados. Logo, o AFD apresentado já
é mínimo.

\subsection{Operadores e delimitadores}

No reconhecimento da união dos símbolos, os estados finais $q_1$, $q_2$ e
$q_3$ têm comportamento equivalente quanto à aceitação, pois não possuem
transições posteriores relevantes. Contudo, no analisador léxico, eles
representam saídas semânticas diferentes: operador aritmético, atribuição e
delimitador. Por essa razão, são mantidos separados para preservar o tipo
de token produzido.

\section{Implementação de apoio em Python}

Além da representação formal, foi elaborada uma implementação de apoio em
Python. O código utiliza expressões regulares para localizar os padrões
léxicos e classificar cada sequência encontrada. As palavras-chave são
identificadas após o reconhecimento de um identificador, verificando se o
texto pertence ao conjunto de palavras reservadas.

\begin{lstlisting}[caption={Implementação simplificada do analisador léxico da MiniLang.},label={lst:lexer}]
import re

PALAVRAS_CHAVE = {"if", "else", "while", "int", "float"}

PADRAO = re.compile(
    r"(?P<ESPACO>\s+)"
    r"|(?P<DECIMAL>[0-9]+\.[0-9]+)"
    r"|(?P<INTEIRO>[0-9]+)"
    r"|(?P<IDENTIFICADOR>[a-zA-Z_][a-zA-Z0-9_]*)"
    r"|(?P<OPERADOR>[+\-*/=])"
    r"|(?P<DELIMITADOR>[(){};])"
)

def analisar(codigo):
    tokens = []
    posicao = 0

    while posicao < len(codigo):
        resultado = PADRAO.match(codigo, posicao)

        if resultado is None:
            caractere = codigo[posicao]
            raise ValueError(
                f"Caractere invalido na posicao {posicao}: {caractere}"
            )

        tipo = resultado.lastgroup
        valor = resultado.group()

        if tipo != "ESPACO":
            if tipo == "IDENTIFICADOR" and valor in PALAVRAS_CHAVE:
                tipo = "PALAVRA_CHAVE"

            tokens.append((tipo, valor))

        posicao = resultado.end()

    return tokens


codigo = "int idade = 20;"
for token in analisar(codigo):
    print(token)
\end{lstlisting}

O exemplo de entrada \texttt{int idade = 20;} contém uma palavra-chave,
um identificador, um símbolo de atribuição, um número inteiro e um
delimitador. A classificação esperada é apresentada na
Tabela~\ref{tab:exemplo}.

\begin{table}[H]
\centering
\caption{Classificação dos tokens da instrução de exemplo}
\label{tab:exemplo}
\begin{tabular}{ll}
\toprule
\textbf{Token} & \textbf{Classificação} \\
\midrule
\texttt{int} & Palavra-chave \\
\texttt{idade} & Identificador \\
\texttt{=} & Atribuição \\
\texttt{20} & Número inteiro \\
\texttt{;} & Delimitador \\
\bottomrule
\end{tabular}
\end{table}

A implementação percorre a entrada da esquerda para a direita. Em cada
posição, tenta reconhecer um dos padrões definidos. Espaços em branco são
ignorados, enquanto os demais elementos são armazenados como pares
\texttt{(tipo, valor)}. Caso nenhum padrão corresponda ao trecho atual, é
gerado um erro informando a posição do caractere inválido.

\section{Análise da implementação}

A implementação utiliza expressões regulares para reconhecer os padrões
léxicos. A busca por cada padrão ocorre a partir da posição atual do texto,
evitando que caracteres inválidos sejam silenciosamente ignorados. A
verificação de palavras reservadas é realizada depois do reconhecimento
genérico de identificadores, pois palavras-chave e identificadores possuem
estrutura textual semelhante.

A função de análise mantém a posição atual da leitura e avança até o final
do código-fonte. Dessa forma, cada trecho reconhecido é consumido uma única
vez. Para uma entrada com $n$ caracteres, a quantidade de posições
percorridas é linear; o custo efetivo também depende da quantidade e da
complexidade dos padrões utilizados pela biblioteca de expressões regulares.

É importante destacar que o código apresentado é uma implementação de apoio
para as categorias especificadas. Ele não constitui um compilador completo:
não realiza análise sintática, verificação semântica ou geração de código.

\section{Conclusão}

O desenvolvimento do analisador léxico para a MiniLang permitiu relacionar
os conceitos de Linguagens Formais e Autômatos com uma aplicação prática
presente no processo de compilação. Foram descritas as categorias léxicas,
as expressões regulares, as definições formais dos AFDs, os diagramas de
transição e as tabelas correspondentes.

A discussão sobre minimização mostrou como estados podem ser comparados
pelo comportamento de suas transições e pela condição de aceitação. Também
foi observado que, em um analisador léxico, a classificação dos tokens deve
ser preservada, mesmo quando estados finais são equivalentes apenas do ponto
de vista da linguagem reconhecida.

Por fim, a implementação em Python demonstrou uma forma de percorrer o
código-fonte, reconhecer padrões e produzir tokens. Como continuidade,
podem ser acrescentados novos operadores, regras léxicas, mensagens de erro
mais detalhadas e integração com uma etapa de análise sintática.

\end{document}
Você atingiu seu limite de análise de dados. Faça upgrade para ChatGPT Plus ou tente novamente amanhã após 18:48.




